\documentclass[12pt]{article}

% هوامش
\usepackage{geometry}
\geometry{margin=0.8in}

% ألوان/رؤوس/تذييل/تباعد
\usepackage{xcolor}
\usepackage{fancyhdr}
\usepackage{setspace}
\usepackage{enumitem}

% دعم يونيكود ولغات (يستلزم XeLaTeX)
\usepackage{fontspec}
\usepackage{polyglossia}

% Math packages for bmatrix/pmatrix/align etc.
\usepackage{amsmath, amssymb, mathtools}

% لصفّ الخيارات أفقياً مع لفّ تلقائي
\usepackage{tabularx, ragged2e, array}
\newcolumntype{Y}{>{\RaggedRight\arraybackslash}X}

\setmainlanguage[numerals=western]{arabic} % أرقام 0-9
\setotherlanguage{english}

% الخطوط (بدّل الاسم إذا لزم)
\newfontfamily\arabicfont[Script=Arabic,Scale=1.05]{Amiri}
\newfontfamily\englishfont{Times New Roman}

% ترويسة/تذييل (نفس التموضع)
\pagestyle{fancy}
\fancyhf{}
\renewcommand{\headrulewidth}{0pt}
\cfoot{\textenglish{Page } \thepage}

% ماكرو الخيار (وسيطان فقط)
\newcommand{\choice}[2]{\textbf{#1}.~#2}

% فقرات
\setlength{\parskip}{0.4em}
\setlength{\parindent}{0pt}

\begin{document}
% ======= Header (fixed layout) =======
\noindent
\begin{tabular*}{\textwidth}{@{\extracolsep{\fill}} l c r}
\textenglish{Date: {\textenglish{September 4, 2025}}} &
{\Large \textbf{{\textenglish{Algorithms -1-}}}} &
{\textenglish{University of Aleppo}} \\
\textenglish{Student Name:} &
\textenglish{Model: {\textenglish{A}}} &
\textenglish{\textbf{Duration:} {\textenglish{90 minuts}}} \\
\textenglish{Student Number:} & & \\
\end{tabular*}

\vspace{0.9cm}

\section*{اختر الإجابة الصحيحة لكل سؤال}
\setstretch{1.25}

\noindent \textbf{\textenglish{1.}} {\textarabic{احسب مشتق التابع }}\( f(x) = x^2\){\textarabic{ وأوجد نهايته عند اللانهاية }}\(\lim_{x \to \infty}f(x)\)

\begin{tabularx}{\linewidth}{@{}*{5}{Y}@{}}
\choice{A}{{{\textenglish{14}}}} & \choice{B}{{{\textenglish{15}}}} & \choice{C}{{{\textenglish{16}}}} & \choice{D}{{{\textenglish{17}}}} & \choice{E}{{{\textenglish{17}}}} \\
\end{tabularx}
\hfill {\textenglish{[1 Point]}}\\[0.5em]

\noindent \textbf{\textenglish{2.}} {\textenglish{Please calculate }}\(\int f(x)dx\)

\begin{tabularx}{\linewidth}{@{}*{4}{Y}@{}}
\choice{A}{{{\textenglish{13}}}} & \choice{B}{{{\textenglish{14}}}} & \choice{C}{{{\textenglish{15}}}} & \choice{D}{{{\textenglish{16}}}} \\
\end{tabularx}
\hfill {\textenglish{[2 Points]}}\\[0.5em]

\noindent \textbf{\textenglish{3.}} {\textenglish{Test}}

\begin{tabularx}{\linewidth}{@{}*{5}{Y}@{}}
\choice{A}{{{\textenglish{13}}}} & \choice{B}{{{\textenglish{14}}}} & \choice{C}{{{\textenglish{15}}}} & \choice{D}{{{\textenglish{16}}}} & \choice{E}{{{\textenglish{17}}}} \\
\end{tabularx}
\hfill {\textenglish{[3 Points]}}\\[0.5em]

\noindent \textbf{\textenglish{4.}} {\textenglish{Test2}}

\begin{tabularx}{\linewidth}{@{}*{5}{Y}@{}}
\choice{A}{{{\textenglish{1}}}} & \choice{B}{{{\textenglish{2}}}} & \choice{C}{{{\textenglish{3}}}} & \choice{D}{{{\textenglish{4}}}} & \choice{E}{{{\textenglish{5}}}} \\
\end{tabularx}
\hfill {\textenglish{[3 Points]}}\\[0.5em]

\noindent \textbf{\textenglish{5.}} {\textenglish{Test3}}

\begin{tabularx}{\linewidth}{@{}*{5}{Y}@{}}
\choice{A}{{{\textenglish{1}}}} & \choice{B}{{{\textenglish{2}}}} & \choice{C}{{{\textenglish{3}}}} & \choice{D}{{{\textenglish{4}}}} & \choice{E}{{{\textenglish{5}}}} \\
\end{tabularx}
\hfill {\textenglish{[1 Point]}}\\[0.5em]

\noindent \textbf{\textenglish{6.}} {\textenglish{test4}}

\begin{tabularx}{\linewidth}{@{}*{5}{Y}@{}}
\choice{A}{{{\textenglish{1}}}} & \choice{B}{{{\textenglish{2}}}} & \choice{C}{{{\textenglish{3}}}} & \choice{D}{{{\textenglish{4}}}} & \choice{E}{{{\textenglish{5}}}} \\
\end{tabularx}
\hfill {\textenglish{[1 Point]}}\\[0.5em]

\noindent \textbf{\textenglish{7.}} {\textenglish{Hello!}}

\begin{tabularx}{\linewidth}{@{}*{4}{Y}@{}}
\choice{A}{{{\textenglish{1}}}} & \choice{B}{{{\textenglish{2}}}} & \choice{C}{{{\textenglish{3}}}} & \choice{D}{{{\textenglish{4}}}} \\
\end{tabularx}
\hfill {\textenglish{[2 Points]}}\\[0.5em]

\noindent \textbf{\textenglish{8.}} {\textenglish{Find the inverse of the following matrix: }}\(\begin{bmatrix} 1 & 3 & 4 \\ 5 & 2 & 1 \\ 6 & 8 & 1 \end{bmatrix}\)

\begin{tabularx}{\linewidth}{@{}*{4}{Y}@{}}
\choice{A}{{{\textenglish{There is no inverse for this matrix}}}} & \choice{B}{{\[\begin{pmatrix} 1 & 2 & 3 \\ 4 & 5 & 6 \\ 7 & 8 & 9  \end{pmatrix}\]}} & \choice{C}{{{\textenglish{0}}}} & \choice{D}{{{\textenglish{0}}}} \\
\end{tabularx}
\hfill {\textenglish{[4 Points]}}\\[0.5em]

\noindent \textbf{\textenglish{9.}} {\textenglish{Solve for }}\(x \){\textenglish{ in the equation }}\(2x+6 = 0 \)

\begin{tabularx}{\linewidth}{@{}*{4}{Y}@{}}
\choice{A}{{{\textenglish{-1}}}} & \choice{B}{{{\textenglish{-2}}}} & \choice{C}{{{\textenglish{-3}}}} & \choice{D}{{{\textenglish{3}}}} \\
\end{tabularx}
\hfill {\textenglish{[1 Point]}}\\[0.5em]

\noindent \textbf{\textenglish{10.}} \(x=3+6 \){\textenglish{ what is the value of x?}}

\begin{tabularx}{\linewidth}{@{}*{5}{Y}@{}}
\choice{A}{{{\textenglish{7}}}} & \choice{B}{{{\textenglish{9}}}} & \choice{C}{{{\textenglish{8}}}} & \choice{D}{{{\textenglish{10}}}} & \choice{E}{{{\textenglish{0}}}} \\
\end{tabularx}
\hfill {\textenglish{[1 Point]}}\\[0.5em]



\end{document}
